\renewcommand{\figurename}{\protect\bf Figure}



\def\stat{semenikhin}



\def\tit{NEURAL NETWORK ARCHITECTURE FOR~ARTIFACTS~DETECTION~IN~ZTF~SURVEY}





\def\titkol{Neural network architecture for artifacts detection in~ZTF survey}



\def\aut{T.\,A.~Semenikhin$^1$}



\def\autkol{T.\,A.~Semenikhin}



\titel{\tit}{\aut}{\autkol}{\titkol}



\renewcommand{\thefootnote}{\arabic{footnote}}

\footnotetext[1]{Sternberg Astronomical Institute, M.\,V.~Lomonosov Moscow State University, 

13~Universitetsky Prosp., Moscow 119234, Russian Federation}





\index{Семенихин Т.\,А.}

\index{Semenikhin T.\,A.} 



\def\leftfootline{\small{\textbf{\thepage}

\hfill Sistemy i Sredstva Informatiki~---

Systems and Means of Informatics 2024 vol~34 no\,1}

}%

 \def\rightfootline{\small{Sistemy i Sredstva Informatiki~---

 Systems and Means of Informatics   2024   vol~34   no\,1

\hfill \textbf{\thepage}}}



      \label{st\stat}



\vspace*{-8pt}



\Abste{Today, astronomers are faced with a challenge of handling vast volume of 

data as modern astronomical surveys are capable of generating terabytes of data 

in a single night. One of such survey is the Zwicky Transient Facility (ZTF), an 

automated sky survey that provides approximately a million alerts per 

observational night. However, a~significant part of the detected objects turn 

out to be artifacts, i.\,e., phenomena of a nonastrophysical origin. Therefore, 

specialists must invest time in manually classifying these objects as there is 

currently no efficient method that can perform this task without human 

intervention. The goal of the work is the development of an algorithm to 

predict whether a light curve from the ZTF data releases (DRs) 

has a~bogus nature or not, based on the sequence of frames. A~labeled dataset 

provided by experts from SNAD team was utilized, comprising 2230 frames 

series. Due to  substantial size of the frame sequences, the application of a 

variational autoencoder (VAE) was deemed necessary for mapping the images into 

lower-dimensional vectors. For the task of binary classification based on sequences of 

compressed frame vectors, a~recurrent neural network (RNN) was employed. Several 

neural network models were considered and the quality metrics were assessed 

using k-fold cross-validation. The final performance metrics, including $ 

\rm{ROC-AUC}=0.86 \pm 0.01$ and $\rm{Accuracy}=0.80 \pm 0.02$, suggest that the 

model has practical utility. The code implementing the algorithm is available on 

{\tt GitHub}.}





\KWE{neural network; data analysis; real-bogus classification}



\DOI{10.14357/08696527240106}{LRMTGD}



      \thispagestyle{myheadings}



\vskip 12pt plus 9pt minus 6pt



\vspace*{-14pt}





\section{Introduction}



\noindent

Zwicky Transient Facility~\cite{2019PASP..131f8003B, 2019PASP..131a8002B} is an automated sky survey. The telescope generates 

a~large amount of data every night. The survey data consist of images capturing 

specific parts of the sky which are later used to create DRs.



The purpose of creating DRs is to investigate any variable sources in 

the sky. The survey frames are processed using the widely-used algorithm 

SExtractor~\cite{1996A&AS..117..393B} which detects sources, determines 

their coordinates, and performs photometry on them. Subsequently, all variable 

sources are assigned a~unique object identifier (OID) which is added to the so-called 

ZTF DRs. It is assumed that there should be no nonperiodic variable 

objects (such as supernovae, red dwarf flares, etc.)\ in the DRs. However, due to 

the fact that the applied algorithms are not perfect, unusual objects are 

encountered in these 

data~\cite{2021MNRAS.502.5147M,2023A&A...672A.111P,Chan_2022}. This motivates 

researchers, even those whose scientific interests are not related to variable 

stars, to work with DRs.



    \begin{figure} %fig1

      \vspace*{1pt}

  \begin{center}

 \mbox{%

 \epsfxsize=109.997mm 

\epsfbox{sem-1.eps}

 }

\end{center}

\vspace*{-6pt}

        \Caption{Examples of artifacts found within 

ZTF data. Images size is $28\times28$~pixels. Frame intensity inverted for better 

readability}

\label{artfs}

                  \end{figure}

                  

 Unfortunately, among the ZTF data, so-called artifacts are often encountered 

(Fig.~\ref{artfs}). An artifact is commonly referred to as any phenomenon of 

nonastrophysical nature. These can be effects related to instrument malfunction 

(defocusing, saturated columns in the CCD array, etc.) or effects associated 

with external conditions (atmospheric turbulence, clouds, bright satellites, 

etc.). Under certain circumstances, these phenomena can be classified as 

astrophysical objects by machine learning methods which significantly 

complicates working with ZTF data.



Classifying astrophysical objects is one of the primary and most challenging 

tasks when working with DR data because they lack class labels. As a result, 

instead of directly using a sequence of object images, photometric observations 

are often employed. Since the observational properties of many astrophysical 

classes have been sufficiently studied, a set of features can be derived from 

the light curve (temporal sequence of photometric observations) that allows for 

their differentiation. Subsequently, classification methods from classical 

machine learning can be applied to the constructed feature sets to solve the 

classification task. However, this approach can lead to situations where 

astrophysical phenomena and artifacts closely resemble each other in their 

photometric representation. For example, a passing bright satellite may 

temporarily illuminate a specific region of the detector which can be mistaken 

for a brief flare of a red dwarf appearing as a short-lived peak in the light 

curve. Therefore, when tackling such problems, specialists need to spend time 

validating the obtained results. Artifacts can account for $68\%$ of the total 

number of objects labeled as astrophysical objects by machine learning 

methods~\cite{2021MNRAS.502.5147M}.



This work is dedicated to implementing an algorithm that would avoid the 

involvement of specialists in classifying artifact/nonartifact objects from DR.

%Currently, there is no efficient approach that solves this problem.

Currently, no attempts have been found in the literature to solve this 

problem.



\section{Data}



\noindent

The labeled object dataset has been provided by SNAD team 

({\sf https://snad.space/}) and does not publicly available. Each object has 

a~unique index, OID, and a~set of labels indicating whether it is an artifact or 

not, along with additional details about the nature of the object (e.\,g., AGN~---

active galactic nucleus, variable, etc. Similarly, for artifacts: ghost, 

defocusing, etc.). For each object, there is a~series of telescope images 

captured at different times. To obtain a~list of links to all available images 

for a~given OID, an API service 

({\sf http://finder.fits.ztf.snad.space/\linebreak  api/docs\#operation/Get\_URLs\_for\_all\_exposures\_by\_object\_ID\_api\_v1\_urls\_by\_oid\_\linebreak get}) is used.







A telescope frame is an image of size $3072\times3080$ depicting a portion of the sky. 

However, for this task, only a~small area within this frame associated with 

a~specific object is of interest. Using OID in the ZTF DR, coordinates are 

determined for the corresponding object (source). In turn, each telescope frame 

is tied to a~coordinate grid. Thus, knowing the OID of the source allows 

determining its center on the telescope frame. To download a specific frame via 

a~link, IRSA service 

({\sf https://irsa.ipac.caltech.edu/docs/program\_interface/ztf\_api.html}) is 

used where, upon request, one can specify the coordinates of the source's center 

and the dimensions of the rectangular area. According to these specifications, 

the desired region will be cropped from the original image. This avoids 

downloading the entire large-sized image allowing the entire dataset to be 

downloaded in approximately one day.





The size of the cropped area was chosen to be $28\times28$~pixels since the 

angular size of the sources of interest does not exceed $\sim$10 arcsec (1 pixel 

in ZTF frames corresponds to~1.012~arcsec). Thus, the region is generously 

sized and the input image size matches that of the popular MNIST~\cite{6296535} 

dataset. If object was close to the image border, the corresponding missing 

areas were filled with zeros (this is necessary for the correct operation of the 

{\tt PyTorch} library). The frames have their background intensity values; so, 

each image needs to be normalized. Normalization corresponding to 

equation

\begin{equation*}

 \label{norm}

    x_{\mathrm{norm}} = \fr{x - \min(x)}{\max(x) - \min(x)}

\end{equation*}



\pagebreak



\



\vspace*{-12pt}



            \begin{floatingfigure}{74mm} %fig2

              \vspace*{-2pt}

  \begin{center}

 \mbox{%

 \epsfxsize=62.893mm 

\epsfbox{sem-2.eps}

 }

\end{center}

\vspace*{-9pt}

                \Caption{Distribution of images number per object}

                \vspace*{6pt}

                \label{hist}

            \end{floatingfigure}

            

            \noindent

where $x$ is the frame; $\mu$ is the mean intensity value of the frame; 

and $\sigma$ is the standard deviation of the frame intensity  was applied. Figure~\ref{artfs} shows examples of preprocessed frames.



The dataset contains 2230 objects including 1150~artifacts and 1080 

astrophysical objects, totaling 1,015,177~images.





One of the main challenges of the task is that objects have a large number of 

observations (Fig.~\ref{hist}) and it is unknown at which specific time 

points anomalies were observed. Therefore, an algorithm capable of working with 

sequences of images rather than individual anomalous images is necessary.





\section{Methods}



\noindent

The implemented approach involves two stages. First, VAE is trained on all images in the dataset to obtain informative compressed 

representations of frames. Then, an RNN is trained on 

sequences of these compressed representations to solve the classification task. 

Combining these two stages into a~single RNN with 

convolutional layers at the beginning is not feasible due to computational 

resource limitations. All neural networks are implemented using the {\tt Python} 

library {\tt PyTorch}.



\subsection{Variational autoencoder}



\noindent

The implemented encoder consists of 5 convolutional layers (with image channel 

dimensions: 32, 64, 128, 256, and 512, respectively) with the LeakyReLU activation 

function and batch normalization after each layer. Each output convolutional 

layer returns half the size of the input image; so, the encoder output is a~vector whose size is a model parameter. The architecture of the decoder is 

symmetrical to that described above, except that transposed convolutional layers 

are used instead of convolutional layers.





The VAE was trained on all images from dataset. When selecting the parameters 

for the VAE, different levels of image compression were considered. In machine 

learning literature, the vector obtained at the output of the encoder is 

commonly referred to as the latent state. In Fig.~\ref{vae_loss}, curves~\textit{1} and~\textit{2}

 represent the loss curves for latent state sizes of~36 and~78, 

respectively. It can be observed that increasing the latent state size requires 

more training epochs. In this work, a~latent state size\linebreak\vspace*{-12pt}



\pagebreak



\



\vspace*{-12pt}





            \begin{floatingfigure}{78mm} %fig3

              \vspace*{-4pt}

  \begin{center}

 \mbox{%

 \epsfxsize=66.873mm 

\epsfbox{sem-3.eps}

 }

\end{center}

\vspace*{-11pt}

                 \Caption{Loss functions during training a~VAE: \textit{1}~--- $\mathrm{latent~size}=36$;

                 \textit{2}~--- $\mathrm{latent~size}\protect\linebreak =78$; and \textit{3}~--- random rotate}

                 \vspace*{9pt}

                  \label{vae_loss}

            \end{floatingfigure}

            



\noindent 

 of~36 proved to be 

optimal. Image augmentation through random rotation was also

explored (curve~\textit{3} 

in Fig.~\ref{vae_loss}); however, it was found that this significantly increased 

the number of required training epochs without a visible improvement in the 

model. Therefore, the final model does not include any augmentations. For model 

training, the Adam optimizer~\cite{kingma2017adam} was used with 

a~$\mathrm{learning~rate} = 5 \cdot 10^{-5}$ for~100~epochs. The core idea of the Adam 

optimizer is to combine adaptive moment estimation with adaptive learning rates 

for efficient model optimization in machine learning. The loss function for the 

VAE takes the following form:

    \begin{multline*}

            L = \mathrm{LogCosh}\left(x, \hat{x}\right) + \omega_{\mathrm{kld}}  \mathrm{KL}\,(N(\mu, \sigma) || 

N(0, 1)) \\

         {}=   \ln{\cosh\left(x - \hat{x}\right)} + \omega_{\mathrm{kld}} \left( \ln 

\fr{1}{\sigma} + \fr{\sigma^2 + \mu^2}{2} - \fr{1}{2} \right).

        \end{multline*}

Here, $x$ and $\hat{x}$ represent the original image and its reconstruction by the 

decoder, respectively; $\mu$ and~$\sigma$ are vectors of means and variances 

obtained from the encoder's output; and $\omega_{\mathrm{kld}} = 8\cdot 10^{-5}$ serves as 

the weight coefficient for the Kullback--Leibler divergence.



By using the compressed representations of the images, it is expected that the 

compressed representations of anomalous frames are outliers from the overall 

distribution of normal images. After the VAE was trained, the obtained 

compressed representations of frames were saved. The next stage of the work 

involved training an RNN for binary classification.



\subsection{Recurrent neural network}



\noindent

The fundamental idea behind RNNs is the utilization of hidden states that 

capture information from previous time steps allowing them to maintain a form 

of memory and learn dependencies within sequential data. By using recurrent 

connections, RNNs can model sequences of arbitrary length making them powerful 

tools for tasks involving sequential information. However, traditional RNNs have 

issues with vanishing and exploding gradients leading to difficulties in 

learning long-range dependencies. To address this, various RNN architectures 

like LSTM (Long Short-Term Memory)~\cite{10.1162/neco.1997.9.8.1735} and GRU 

(Gated Recurrent Unit)~\cite{cho2014learning} have been developed which 

introduce gating mechanisms to better control the flow of information and 

mitigate gradient-related problems.



The model takes inputs of size (batch~size, sequence~length, and latent\_dim) where 

latent\_dim represents the dimensionality of the compressed representation. 

Embeddings for object frames along with their corresponding class labels are 

fed into an RNN. The model parameters are optimized using the Adam optimizer with 

a~$\mathrm{learning~rate}=10^{-4}$ for~500~epochs.







\begin{table}[b]\small

\begin{center}

\begin{tabular}{lccc}

\multicolumn{4}{p{100mm}}{Model results. The metric values are averaged over 5 test folds as 

well as the standard deviation}\\

\multicolumn{4}{p{100mm}}{\ }\\[-6pt]

\hline

\multicolumn{1}{c}{Model name} & ROC-AUC & Accuracy & \multicolumn{1}{c}{F1-score} \\

\hline

&&\\[-10pt]

GRU & $0.84 \pm 0.01$ & $0.78 \pm 0.02$ & $0.79 \pm 0.01$\\

Bidirectional GRU\;+\;$L_2$ & $0.86 \pm 0.01$ & $0.80 \pm 0.02$ & $0.80 \pm 0.01$\\

GRU\;$+$\;Tversky loss & $0.84 \pm 0.02$ &  $0.79 \pm 0.02$ & $0.79 \pm 0.02$\\

LSTM & $0.76 \pm 0.04$ & $0.76 \pm 0.03$ & $0.77 \pm 0.02$ \\

\hline

\end{tabular}

\end{center}

\end{table}











The baseline model considered was a GRU layer with a hidden state size of~128 

and the loss function used during training the RNN was 

cross-entropy. In addition to the baseline model, the following variations were 

explored: a~bidirectional GRU cell with $L_2$ weight regularization 

($\mathrm{weight~decay}=10^{-5}$); a~GRU layer with a~modified loss function (the Tversky 

Index~\cite{salehi2017tversky}) term was added to the loss function); and an LSTM 

layer with the same parameters as the baseline model.



During the RNN training, we employed k-fold cross-validation ($k=5$). The 

dataset was randomly divided into 5~folds and then 5 models were trained, each 

having a~specific fold as the test set and the union of the remaining~4~folds as 

the training set. Thus, the data split into training\!/\!test sets was unique for 

each model. The final quality metrics for each considered model were calculated 

as the average values across 5 models from the k-fold cross-validation split.



Applying the obtained model to the entire ZTF DR dataset is 

impractical as downloading all images for all objects in the ZTF DR would take 

months or even years. In practice, it is expected that the classifier's decision 

function will serve as an additional source of information for the active 

(employing partial teacher involvement through feedback loops) anomaly detection 

algorithm and the classifier will be applied only to a small subset of objects~--- candidates for anomalies~\cite{2021A&A...650A.195I}. In this case, 

determining the specific threshold of the decision function is the prerogative 

of the active anomaly detection algorithm based on expert feedback. Moreover, 

technically, such a threshold may vary for different feature space regions. 

Therefore, ROC-AUC was chosen as the main metric for evaluating model 

performance.







Bidirectional GRU is different from the regular GRU in that it processes input 

sequences in both forward and backward directions which enables it to capture 

contextual information from past and future time steps. However, when using\linebreak\vspace*{-12pt}



\pagebreak



\



\vspace*{-12pt}



            \begin{floatingfigure}{70mm} %fig4

              \vspace*{-5pt}

  \begin{center}

 \mbox{%

 \epsfxsize=62.717mm 

\epsfbox{sem-4.eps}

 }

\end{center}

\vspace*{-11pt}

                \Caption{The ROC curves examples of trained models for 

one of the cross-validation splits: \textit{1}~--- GRU; \textit{2}~--- LSTM; \textit{3}~--- bidirectional GRU\;+\;$L_2$; and \textit{4}~--- GRU\;+\;Tversky loss

}

%\vspace*{6pt}

\label{roc}

            \end{floatingfigure}

            

            \noindent

             such 

an approach in this task, the model experienced overfitting; so, $L_2$ weight 

regularization was added to it. The table %~\ref{table} 

as well as 

Fig.~\ref{roc} show the main results

             for the considered models. The table 

indicates that LSTM performs worse than the other models. The three remaining 

models exhibit approximately the same level of quality based on the performance 

metrics. However, the bidirectional GRU\;+\;$L_2$ appears slightly more stable 

across different data splits. Therefore, this model was selected for further 

work.







\section{Concluding Remarks}



\noindent

Throughout the course of this study, sequences of ZTF object frames from 

a~labeled dataset were downloaded and preprocessed. Using the {\tt PyTorch} 

library, a~VAE was implemented and trained enabling the 

compression of images into informative lower-dimensional vectors. Several models 

of RNNs were considered for the task of binary 

classification based on sequences of compressed image vectors. The quality 

metrics of the models were evaluated using k-fold cross-validation, with the 

best-performing model achieving an $ \mathrm{ROC}\mbox{-}\mathrm{AUC}=0.86 \pm 0.01$ and 

$\mathrm{Accuracy}=0.80 \pm 0.02$. All the code used in the study is available on 

{\tt GitHub} ({\sf https://github.com/semtim/RB\_ZTF}).



Thus, an algorithm has been developed that aims to classify objects from ZTF DR 

in a manner similar to that of a specialist who annotated the training dataset. 

This algorithm stands out from others by utilizing sequences of observation 

frames rather than photometric time series. While computationally more 

challenging, images contain more information about the object's radiation than 

its corresponding photometry.





The next stage of the work will involve: searching for the optimal architecture 

for both the autoencoder and the recurrent network; selecting the most effective 

loss function for the classifier in terms of our task; analyzing the dependence 

of the first and second type errors on the chosen threshold value; and conducting 

computational experiments using different optimizers. This algorithm can be 

employed within the scope of the SNAD project as 

the obtained quality metrics indicate its practical value. Since the training 

dataset used in this work was annotated by specialists from the SNAD project, it 

is expected that the implementation of this algorithm into the general pipeline 

will reduce the number of artifacts among anomalies. This will enable 

specialists to detect more new astrophysical objects.











\Ack



\vspace*{-3pt}



\noindent

The paper is published on the proposal of the Program Committee of the 25th 

International Conference on Data Analytics and Management in Data Intensive 

Domains (DAMDID/RCDL 2023). The author thanks the SNAD team for the provided 

comments and support. The author acknowledges the support of M.\,V.~Lomonosov 

Moscow State University Program of Development. The work was supported by Nonprofit 

Foundation for the Development of Science and Education ``Intellect.''



\renewcommand{\bibname}{\protect\rmfamily References}



{\small\frenchspacing

{%\baselineskip=10.8pt

\begin{thebibliography}{99}





%1

\bibitem{2019PASP..131a8002B}

\Aue{Bellm, E.\,C., S.\,R.~Kulkarni, M.\,J.~Graham, \textit{et al.}} 2019. The Zwicky 

Transient Facility: System overview, performance, and first results. 

\textit{Publ. Astron. Soc. Pac.} 131(995):018002. 19 p. doi: 10.1088\!/\!1538-3873\!/\!aaecbe.



%2

\bibitem{2019PASP..131f8003B}

\Aue{Bellm, E.\,C., S.\,R.~Kulkarni, T.~Barlow, \textit{et al.}} 2019.

The Zwicky Transient Facility: Surveys and scheduler. \textit{Publ. Astron. Soc. Pac.} 131(1000):068003. 13~p. 

doi: 10.1088\!/\!1538-3873/ab0c2a.



%3

\bibitem{1996A&AS..117..393B}

\Aue{Bertin, E., and S.~Arnouts}. 1996. SExtractor: Software for source extraction. 

\textit{Astron. Astrophys. Sup.} 117(2):393--404. doi: 10.1051\!/\!aas:1996164.



%4

\bibitem{2021MNRAS.502.5147M}

\Aue{Malanchev, K.\,L., M.\,V.~Pruzhinskaya, V.\,S.~Korolev, \textit{et al.}} 2021. 

{Anomaly detection in the Zwicky transient facility DR3}. \textit{Mon. Not. R. 

Astron. Soc.} 502(4):5147--5175. doi: 10.1093\!/\!mnras\!/\!stab316.



%5

\bibitem{Chan_2022}

\Aue{Ho-Sang, C., V.\,A.~Villar, S.~Cheung, \textit{et al.}} 2022. Searching for 

anomalies in the ZTF catalog of periodic variable stars. \textit{Astrophys.~J.} 

932(2):118. 21~p. doi: 10.3847\!/\!1538-4357\!/\!ac69d4.



%6

\bibitem{2023A&A...672A.111P}

\Aue{Pruzhinskaya, M.\,V., E.\,E.\,O.~Ishida, A.\,K.~Novinskaya, \textit{et al.}}  2023. 

Supernova search with active learning in ZTF DR3. \textit{Astron. Astrophys.} 

672:A111. 22~p. doi: 10.1051\!/\!0004-6361\!/\!202245172.







%7

\bibitem{6296535}

Deng, L. 2012. The MNIST database of handwritten digit images for machine 

learning research. \textit{IEEE Signal Proc. Mag.} 29(6):141--142. doi: 

10.1109\!/\!MSP.2012.2211477.



%8

\bibitem{kingma2017adam}

\Aue{Kingma, D.\,P., and J.~Ba.} 2017. Adam: A~method for stochastic optimization. 

\textit{arXiv.org}. 15 p. Available at: {\sf https://arxiv.org/abs/1412.6980} 

(accessed March~11, 2024).



%9

\bibitem{10.1162/neco.1997.9.8.1735}

\Aue{Hochreiter, S., and J.~Schmidhuber.} 1997. Long short-term memory. \textit{Neural 

Comput.} 9(8):1735--1780. doi: 10.1162\!/\!neco.1997.9.8.1735.



%10

\bibitem{cho2014learning}

\Aue{Cho, K., B.~Merri\"{e}nboer, C.~Gulcehre, \textit{et al.}} 2014. Learning phrase 

representations using RNN encoder-decoder for statistical machine translation. 

\textit{arXiv.org}. 15~p. Available at: {\sf https://arxiv.org/abs/1406.1078} 

(accessed March~11, 2024).



%11

\bibitem{salehi2017tversky}

\Aue{Salehi, S., D.~Erdogmus, and A.~Gholipour.} 2017. Tversky loss function for image 

segmentation using 3D fully convolutional deep networks. \textit{arXiv.org}. 9~p. Available at: 

{\sf https://arxiv.org/abs/1706.05721} (accessed March~11, 2024).



%12

\bibitem{2021A&A...650A.195I}

\Aue{Ishida, E.\,E.\,O.,  M.\,V.~Kornilov, K.\,L.~Malanchev, \textit{et al.}} 2021. Active 

anomaly detection for time-domain discoveries. \textit{Astron. Astrophys.} 

650:A195. 9~p. doi: 10.1051\!/\!0004-6361\!/\!202037709.



\end{thebibliography} } }







\vspace*{-6pt}



\hfill{\small\textit{Received December 6, 2023}}



\vspace*{-6pt}





\Contrl



\noindent

\textbf{Semenikhin Timofey A.} (b.\ 1999)~--- engineer, Sternberg Astronomical 

Institute, M.\,V.~Lomonosov Moscow State University, 13~Universitetsky Prosp., 

Moscow 119234, Russian Federation; \mbox{ofmafowo@gmail.com}













\vspace*{12pt}



\hrule



\vspace*{2pt}



\hrule



\vspace*{6pt}



%\newpage







\def\tit{ПОИСК АРТЕФАКТОВ НА ИЗОБРАЖЕНИЯХ ОБЗОРА ZTF ПРИ~ПОМОЩИ НЕЙРОННЫХ 

СЕТЕЙ$^*$}



\def\aut{Т.\,А.~Семенихин}





\def\titkol{Поиск артефактов на изображениях обзора ZTF при помощи нейронных 

сетей}



\def\autkol{Т.\,А.~Семенихин}



\titel{\tit}{\aut}{\autkol}{\titkol}



\vspace*{-9pt}



\noindent

Государственный астрономический институт имени П.\,К.~Штернберга

Московского государственного университета имени М.\,В. Ломоносова 

% \mbox{ofmafowo@gmail.com}



\vspace*{6pt}



\def\leftfootline{\small{\textbf{\thepage}

\hfill СИСТЕМЫ И СРЕДСТВА ИНФОРМАТИКИ\ \ \ том~34\ \ \ номер~1\ \ \ 2024}

}%

 \def\rightfootline{\small{СИСТЕМЫ И СРЕДСТВА ИНФОРМАТИКИ\ \ \ том~34\ \ \ номер~1\ \ \ 2024

\hfill \textbf{\thepage}}}













{\renewcommand{\thefootnote}{\fnsymbol{footnote}} \footnotetext[1]

{Статья публикуется по представлению программного комитета XXV Международной 

конференции <<Аналитика и управление данными в~областях с~интенсивным 

использованием данных>> (DAMDID\!/\!RCDL 2023). Автор выражает благодарность команде 

SNAD за комментарии и~поддержку. Работа выполнена с~использованием оборудования, 

приобретенного за счет средств Программы развития Московского университета. 

Исследование выполнено при финансовой поддержке Некоммерческого фонда развития 

науки и~образования <<Интеллект>>.}}







\Abst{Сегодня астрономам приходится работать с большими объемами данных, 

поскольку современные инструменты способны генерировать терабайты информации за 

одну ночь. Одним из таких инструментов стал автоматизированный обзор неба Zwicky 

Transient Facility (ZTF), который за одну ночь детектирует порядка миллиона новых 

вспышек. Однако значительная доля найденных объектов оказывается артефактами, т.\,е.\

 явлениями, имеющими неастрофизическую природу. Поэтому специалистам 

приходится тратить время на классификацию объектов вручную, так как на текущий 

момент не существует эффективного метода, который делал бы это без учас\-тия 

человека. Целью данной работы ставилась реализация эффективного алгоритма, который 

по последовательности кадров объекта из обзора ZTF 

предсказывал бы, относится он к~артефактам или нет. Для реализации алгоритма 

использована выборка, размеченная специалистами коллаборации SNAD 

и~содержащая~2230~серий кадров объектов. Так как последовательности кадров 

достаточно велики, использован вариационный автоэнкодер, который позволяет 

отобразить изображение в~вектор меньшей длины. Для решения задачи бинарной 

классификации по последовательности сжатых в~векторы кадров применялась 

рекуррентная нейронная сеть. Были рассмотрены несколько моделей нейронных сетей, 

для оценки метрик качества использовалась k-fold кросс-валидация. Итоговые 

метрики качества составляют $ \mathrm{ROC}\mbox{-}\mathrm{AUC}\hm= 0{,}86 \pm 0{,}01$ и~точ\-ность~$0{,}80 

\pm 0{,}02$ и~позволяют говорить, что модель имеет практическую ценность. Код с 

реализацией алгоритма доступен на %{\sf https://github.com/semtim/RB\_ZTF}

{\tt 

GitHub}.}



\KW{нейронные сети; анализ данных; классификация ре\-аль\-ных/лож\-ных объектов}





\DOI{10.14357/08696527240106}{LRMTGD}



%\vspace*{6pt}





\renewcommand{\bibname}{\protect\rmfamily Литература}

%\renewcommand{\bibname}{\large\protect\rm References}



{\small\frenchspacing

{%\baselineskip=10.8pt

\begin{thebibliography}{99}







\bibitem{2019PASP..131a8002B-1} %1

\Au{Bellm~E.\,C., Kulkarni~S.\,R., Graham~M.\,J., \textit{et al.}}

The Zwicky Transient Facility: System overview, performance, and first results~/\!\!/ Publ. Astron. Soc. Pac., 

2019. Vol.~131. Iss.~995. Art.~018002. 19~p. doi: 10.1088\!/\!1538-3873/aaecbe.



\bibitem{2019PASP..131f8003B-1} %2

\Au{Bellm~E.\,C., Kulkarni~S.\,R., Barlow~T., \textit{et al}}.

{The Zwicky Transient Facility: Surveys and scheduler}~/\!\!/ Publ. Astron. Soc. Pac., 2019. Vol.~131. 

Iss.~1000. Art.~068003. 13~p. doi: 10.1088\!/\!1538-3873/ab0c2a.





\bibitem{1996A&AS..117..393B-1} %3

\Au{Bertin~E., Arnouts~S.} SExtractor: Software for source extraction~/\!\!/ Astron. Astrophys. Sup., 

1996. Vol.~117. No.~2. P.~393--404. doi: {10.1051\!/\!aas:1996164}.



\bibitem{2021MNRAS.502.5147M-1} %4

\Au{Malanchev K.\,L., Pruzhinskaya~M.\,V., Korolev~V.\,S., \textit{et. al.}} Anomaly 

detection in the Zwicky Transient Facility DR3~/\!\!/ Mon. Not. R. 

Astron. Soc., 2021. Vol.~502. 

No.\,4. P.~5147--5175. doi: 10.1093\!/\!mnras/stab316.







\bibitem{Chan_2022-1} %5

\Au{Ho-Sang C., Villar~V.\,A., Cheung~S., \textit{et al.}} Searching for anomalies in 

the ZTF catalog of periodic variable stars~/\!\!/ Astrophys.~J., 2022. Vol.~932. No.\,2. 

Art.~118. 21~p. doi: 10.3847\!/\!1538-4357\!/\!ac69d4.



\bibitem{2023A&A...672A.111P-1} %6

\Au{Pruzhinskaya M.\,V., Ishida~E.\,E.\,O.,  Novinskaya~A.\,K., \textit{et al.}} 

{Supernova search with active learning in ZTF DR3} // Astron. Astrophys., 2023. 

Vol.~672. Art.~A111. 22~p. doi: 10.1051\!/\!0004-6361\!/\!202245172.





\bibitem{6296535-1}

\Au{Deng Li.} The MNIST database of handwritten digit images for machine 

learning research~/\!\!/ IEEE Signal Proc. Mag., 2012. Vol.~29. No.\,6. P.~141--142. doi: 10.1109\!/\!MSP.2012.2211477.



\bibitem{kingma2017adam-1}

 \Au{Kingma D.\,P., Ba~J.} Adam: A~method for stochastic optimization.~--- 

Cornell University, 2017. arXiv:1412.6980 [cs.LG]. 15~p.



\bibitem{10.1162/neco.1997.9.8.1735-1}

\Au{Sepp H., Jurgen~S.} Long short-term memory~/\!\!/ Neural Comput., 1997. 

Vol.~9. No.\,8. P.~1735--1780. doi: 10.1162\!/\!neco.1997.9.8.1735.



\bibitem{cho2014learning-1}

\Au{Cho~K., Merri\"{e}nboer~B., Gulcehre~C., \textit{et al.}} Learning phrase representations using RNN encoder-decoder for 

statistical machine translation.~--- Cornell University, 2014. arXiv:1406.1078 

[cs.CL]. 15~p.



\bibitem{salehi2017tversky-1}

 \Au{Salehi S., Erdogmus~D., Gholipour~A.} Tversky loss function for image 

segmentation using 3D fully convolutional deep networks.~--- Cornell 

University,  2017. arXiv:1706.05721 [cs.CV]. 9~p.



\bibitem{2021A&A...650A.195I-1}

\Au{Ishida~E.\,E.\,O., Kornilov~M.\,V.,  Malanchev~K.\,L., \textit{et al.}} Active 

anomaly detection for time-domain discoveries~/\!\!/ Astron. Astrophys., 2021. 

Vol.~650. Art.~A195. 9~p. doi: 10.1051\!/\!0004-6361\!/\!202037709.





\label{end\stat}



\end{thebibliography}

} }





\hfill{\small\textit{Поступила в редакцию 06.12.2023}}

%\renewcommand{\bibname}{\protect\rm Литература}

\renewcommand{\figurename}{\protect\bf Рис.}

